\chapter[Vanishing criteria, coherence conditions]
{Vanishing criteria, coherence conditions for the sheaves $\SheafExt^{i}_{Y}(F, G)$}\label{VII}

\renewcommand{\theequation}{\thesection.\arabic{equation}}

\section{Study for $i < n$}

Let us prove a lemma:

\begin{lemma} \label{VII.1.1}
Let $X$ be a locally noetherian prescheme, $Y$ a closed subset of~$X$ and $G$ a quasi-coherent $\OX$-Module. Suppose that for every coherent $\OX$-Module~$F$ of support contained in $Y$, one has:
\[
\SheafExt^{n-1}(F, G) = 0.
\]
Then for every coherent $\OX$-Module $F$ and every closed subset $Z$ of $X$ such that $Y \supset \Supp F \cap Z$, one has
\[
\SheafExt^n_Z(F, G) \approx \SheafHom (F, \HH^n_Y(G)).
\]
\end{lemma}

One first notes that
\[
\SheafExt^i_Z(F, G) = \SheafExt^i_{Z \cap \Supp F}(F, G).
\]
(trivial, \cf Expos\'e \Ref{VI}). One first gives the proof for $Z = X$ hence \hbox{$\Supp F \subset Y$}. The functor
\[F \mto \SheafExt^n(F, G),
\]
defined on the category of coherent $\OX$-Modules of support contained in $Y$, is left exact. By virtue of (\Ref{IV}~\Ref{IV.1.3}), it is representable by
\[I = \varinjlim_{k} \SheafExt^n (\OX/\mathcal{I}^{k+1}, G),
\]
where $\mathcal{I}$ is the ideal of definition of $Y$. Now, by (\Ref{II}~\Ref{II.6}), one knows that:
\[
\HH^n_Y(G) \approx \varinjlim_{k} \SheafExt^n (\OX/\mathcal{I}^{k+1}, G).
\]
Whence the conclusion if $ Z= X$. Still by (\Ref{VI}~\Ref{VI.2.3}) one knows that:
\[
\SheafExt^n_Z (F, G) \approx \varinjlim_{k} \SheafExt^n (F/\Jj^{k + 1}F, G),
\]
where $\Jj$ is the ideal of definition of $Z$. The support of $F/\Jj^{k+1}F$ is contained in $Y$ if $Z \cap \Supp F \subset Y$; by what we have just proved, one thus has:
\[
\SheafExt^n_Z (F, G) \approx \varinjlim_{k} \SheafHom (F/\Jj^{k + 1}F, \HH^n_Y(G)).
\]
It remains to show that the natural homomorphism:
\[
\varinjlim_{k} \SheafHom (F/\Jj^{k + 1}F, \HH^n_Y(G)) \to \SheafHom (F, \HH^n_Y (G)),
\]
is an isomorphism when $Z \cap \Supp F \subset Y$. Now $X$ can be covered by affine noetherian open sets; one is thus reduced to the case where $X$ is affine noetherian; then $F(X)$ is an $\OX(X)$-Module of finite type and $\Supp \HH^n_Y(G) \subset Y$. Hence every homomorphism \hbox{$u: F(X) \to \HH^n_Y (G) (X)$} is annihilated by a power of $\Ii$ hence by a power of~$\Jj$, \hfill qed.

\begin{proposition} \label{VII.1.2}
Let $X$ be a locally noetherian prescheme, $Y$ a closed subset of $X$, $G$ a quasi-coherent $\OX$-Module and $n$ an integer. Whatever be $Z$ and $S$, closed subsets of $X$ such that $Z \cap S = Y$, the following conditions are equivalent:
\begin{enumeratei}
\item
$\HH^i_Y (G) = 0$ if $i<n$;

\item
there exists a coherent $\OX$-Module $F$, of support $S$, such that:
\[
\SheafExt^i_Z (F, G) = 0 \text{ if } i<n;\]

\item
for every $\OX$-Module $F$ coherent of support contained in $S$ (\ie $\Supp F \cap Z$ $=$ $\Supp F \cap Y)$, one has:
\[
\SheafExt^i_{Z}(F, G) = 0 \text{ if } i < n;
\]
\item
for every coherent $\OX$-Module $F$, one has:
\[
\SheafExt^i_Y (F, G) = 0 \text{ if } i < n.
\]
\end{enumeratei}
\noindent
Moreover, if they are satisfied, then for every coherent $\OX$-Module $F$ and every closed subset $Z'$ of $X$ such that $Z' \cap \Supp F = Y \cap \Supp F$, one has isomorphisms:
\[
\SheafExt^n_{Z} (F, G) \approx \SheafExt^n_Y (F, G) \approx \SheafHom (F, \HH^n_Y (G)).
\]
\end{proposition}

\begin{proof}
Let us argue by induction. The proposition is trivial for $n < 0$. Suppose it has been proved for $n < q$. If one of the conditions is satisfied for $n = q$, and for two subsets $Z$ and $S$ as stated, by the induction hypothesis one has, for every closed subset $Z'$ of $X$ and every coherent $\OX$-Module $F$ such that $Z' \cap \Supp F = Y \cap \Supp F$, isomorphisms:
\begin{equation}\label{eq:VII.1.1}\SheafExt^{q-1}_{Z'} (F, G) \approx \SheafHom (F, \HH^{q-1}_Y (G)) \approx \SheafExt^{q-1}_Y (F, G).\end{equation}

\noindent Hence:

(i) \ALORS (iv), for one takes $Z' = Y$ in (1.1);

(iv) \ALORS (iii), for one takes $Z' = Z$ in (1.1);

(iii) \ALORS (ii), for one takes $F = \Oo_S$;

(ii) \ALORS (i), for one takes $Z' = Z$ in (1.1); whence $\SheafHom (F, \HH_Y^{q-1} (G)) = 0$; one then notes that:
\[
\Supp \HH_Y^{q-1} (G) \subset Y = Z \cap S \subset S = \Supp F,
\]
and one applies the following lemma:

\begin{lemma} \label{VII.1.3}
Let $X$ be a prescheme, let $P$ be a coherent $\OX$-Module, and let $H$ be a quasi-coherent $\OX$-Module such that:
\[
\SheafHom (P, H) = 0 \text{ and } \Supp P \supset \Supp H.
\]
Then $H = 0$.
\end{lemma}

It suffices to prove the lemma when $X$ is affine, for the affine open sets form a basis of the topology of $X$ and the hypotheses are preserved by restriction to an open set. Now, in this case, one is reduced to a problem on $A$-modules, where $X = \Spec(A)$. One applies the formula (valid under the sole hypothesis that $M$ is of finite type):
\[
\Ass \Hom_A (P, H) = \Supp P \cap \Ass H;\]
One knows that $\Ass H \subset \Supp H \subset \Supp P$ and that $\Ass \Hom_A (P, H) = \emptyset$; hence $\Ass H = \emptyset$, hence $H = 0$.

To finish the proof of the proposition, it remains to note that (iv) allows one to apply~\Ref{VII.1.1}.
\skipqed
\end{proof}

\begin{corollary}\label{VII.1.4}
Let $G$ be a coherent and Cohen--Macaulay $\OX$-Module, let $n \in \ZZ$. The conditions of \Ref{VII.1.2} are equivalent to:
\begin{enumerate}
\item[\textup{(v)}]\hfill$\codim(Y \cap \Supp G, \Supp G) \geq n$.\hfill\null
\end{enumerate}
\end{corollary}
\setcounter{equation}{1}
We first recall that an $\OX$-module is said to be Cohen--Macaulay if, for every $x \in X$, the fibre $G_x$ is a Cohen--Macaulay $\Oo_{X, x}$-module, \ie one has for every $x \in S = \Supp G$:
\begin{equation}\label{eq:VII.1.2}
\prof G_x = \dim G_x = \dim \Oo_{S, x}.
\end{equation}
By Proposition~\Ref{III}~\Ref{III.3.3}, condition (i) of \Ref{VII.1.2} is equivalent to:
\begin{equation}\label{eq:VII.1.3} \prof_Y G = \inf_{x\in Y} \prof G_x \geq n,
\end{equation}
hence also to:
\[
\prof_Y G = \inf_{x \in Y \cap S} \prof G_x \geq n;
\]
for the depth of a zero module is infinite.

Now, by definition:
\[
\codim (Y \cap S, S) = \inf_{x \in S \cap Y} \dim \Oo_{S, x},
\]
whence the conclusion, applying the formula \eqref{eq:VII.1.2}).

We shall now prove a result allowing one to deduce the coherence conditions that we have in view from certain vanishing criteria.

\begin{lemma}\label{VII.1.5}
Let $X$ be a locally noetherian prescheme. Let $T^{\ast}$ be an exact contravariant $\partial$-functor, defined on the category of coherent $\OX$-Modules, with values in the category of $\OX$-Modules. Let $Y$ be a closed subset of $X$. Let $i \in \ZZ$. Suppose that, for every coherent $\OX$-Module of support contained in $Y$, $T^i F$ and $T^{i-1} F$ are coherent. Let $F$ be a coherent $\OX$-Module. For $T^i F$ to be coherent, it is necessary and sufficient that $T^i F''$ be so, where one has set:
\[
F'' = F / \Gamma_Y (F).
\]
\end{lemma}

Indeed, $F' = \Gamma_Y (F)$ is coherent because $X$ is locally noetherian; the exact cohomology sequence of $T^{\ast}$ then gives:
\[
T^{i-1}F' \to T^i F'' \to T^i F \to T^i F'
\]
where the extreme terms are coherent, whence the conclusion.

\begin{lemma}\label{VII.1.6}
If $F$ and $G$ are coherent, and if $\Supp F \subset Y$, $\SheafExt^i_Y (F, G)$ is coherent.
\end{lemma}

Indeed, $\SheafExt^i_Y (F, G)$ is isomorphic to $\SheafExt^i (F, G)$; this is valid on any ringed space $X$, moreover: if $Z$ is a closed subset which contains $Y \cap \Supp F$, $\SheafExt^i_Z (F, G)$ is isomorphic to $\SheafExt^i_Y (F, G)$ (\cf Expos\'e \Ref{VI}).

\begin{proposition}\label{VII.1.7}
\hspace*{3mm}Suppose $F$ and $G$ are coherent and set $\Supp F = S$, $S' = \overline{S \cap (X-Y)}$. Suppose that, for every $x \in Y \cap S'$, one has $\prof G_x \geq n$, then $\SheafExt^i_Y (F, G)$ is coherent for $i < n$.
\end{proposition}

Indeed, \Ref{VII.1.6} allows one to apply \Ref{VII.1.5} to $T^{\ast}(F) = \SheafExt^{\ast}_Y (F, G)$. Setting $F'' = F/ \Gamma_Y (F)$, one sees that $\Supp F'' = S'$. Now, by~\Ref{III}~\Ref{III.3.3}, the hypothesis on the depth of $G$ ensures the vanishing of $\H_{Y\cap S'}(G)$ for $i<n$; by~\Ref{VII.1.2}, one deduces from this the vanishing of $T^i F''$ for $i < n$, whence the conclusion thanks to~\Ref{VII.1.5}.

\section{Study for $i>n$}\label{VII.2}

Let $X$ be a locally noetherian \emph{regular} prescheme, that is, one all of whose local rings are regular. Let $Y$ be a closed subset of $X$. Let $F$ and $G$ be two coherent $\OX$-Modules. Set $S = \Supp F$, $S' = \overline{S \cap (X - Y)}$. Set:
\begin{align*}
m& = \sup_{x \in Y \cap S} \dim \Oo_{X, x},\\
n &= \sup_{x \in Y \cap S'} \dim \Oo_{X, x};
\end{align*}
one has $n \leq m$.

\begin{proposition}\label{VII.2.1}
In the situation described above, one has:
\begin{enumerate}
\item
$\SheafExt^i_Y (F, G) = 0$ if $i > m$,

\item
$\SheafExt^i_Y (F, G)$ is coherent if $i > n$.
\end{enumerate}
\end{proposition}

One first notes that $\SheafExt^i_Y (F, G)$ is coherent for every $i$ when $\Supp F \subset Y$. Moreover, setting as above $F'' = F/\Gamma_Y(F)$, one sees that $\Supp F'' = S'$, hence (2) follows from (1) and from \Ref{VII.1.3}.

To prove (1), one first notes that
\[
\SheafExt^i_Y (F, G) \approx \varinjlim_{k} \SheafExt^i (F/\Jj^k F, G),
\]
where $\Jj$ is the ideal of definition of $Y$. Moreover, it
follows from theorem 4.2.2. of ( A. Grothendieck,
``Sur quelques points d'alg\`ebre
homologique'', \textit{T\^{o}hoku Mathematical Journal}
\textbf{9}  (1957), p. 119--221.) that the $\SheafExt$ commute with
the formation of fibres, at least when $X$ is a locally noetherian
prescheme and when the first argument is coherent;
as the same holds for inductive limits, one finds
isomorphisms:
\[
(\SheafExt^i_Y (F, G))_x \approx \varinjlim_{k} \Ext_{\Oo_{X, x}} ((F/\Jj^k F)_x, G_x)
\]
for every $x \in X$. As $\Supp \Ext^i_Y (F, G) \subset S \cap Y$, it suffices, in order to conclude, to note that $x \in Y \cap S$ implies $\dim \Oo_{X, x} \leq m$, hence:
\[
\Ext^i_{\Oo_{X, x}} ((F/\Jj^k F)_x, G_x) = 0 \text { if } i > m,
\]
for the global cohomological dimension of a regular local ring is equal to its dimension\sfootnote{\Cf $\EGA 0_{\textup{IV}}$ 17.3.1}.

Let $X$ be a locally noetherian prescheme; for every subset $P$ of $X$, set:
\[
D(P) = \left\{ \dim \Oo_{X, p} \mid p \in P \right\}.
\]

\begin{lemma}\label{VII.2.2}
If $P$ is the underlying space of a connected subprescheme of $A$, $D(P)$ is an interval.
\end{lemma}

Indeed, let $a$ and $B$ belong to $D(P)$, corresponding to points $p$ and $q$ of $P$. Let us show that there exists a sequence of points of $P$: $(p = p_1, \dots, p_n = q)$ such that, for $1 \leq i < n $, one has $\left| \dim \Oo_{X, p_i} - \dim \Oo_{X, p_{i+1}} \right| = 1$; it will follow that D(P) contains the interval $[p, q]$. For this, one notes that $p$ and $q$ can be joined by a sequence of irreducible components of $P$, such that two successive components meet. One is reduced to the case where $p$ is the generic point of an irreducible component $Q$ of~$P$, and where $q \in Q$ and hence $q \supset p$, as ideals of $\Oo_q$, where it is trivial from the definition of dimension.

\begin{proposition}\label{VII.2.3} Let $X$ be a locally noetherian regular prescheme, $Y$ a closed subset of $X$ and $F$ a coherent $\OX$-Module. Let $P = Y \cap \Supp F \cap (X - Y)$. Let $n \in \ZZ$, and suppose that $n \not\in D(P)$. Then $\SheafExt^n_Y (F, \OX)$ is coherent.
\end{proposition}

The conclusion is local and the hypotheses are preserved by restriction to an open set. Now $P$ is closed hence locally noetherian, hence locally connected; one may thus assume $X$ affine and noetherian, and $P$ connected. Set $D(P) = [a, b[$, which is legitimate by the preceding lemma; if $n> b$, one concludes by \Ref{VII.2.1}; if $n < a$, one has \hbox{$n < \dim \Oo_{X, x} = \prof \Oo_{X, x}$} for every $x \in P$, and one concludes by \Ref{VII.1.7}.
